% !TEX root = ../main.tex

\chapter{总结与展望}

\section{主要工作与结论}
  本文面向临床实际存在的专家标注成本高、跨中心MRI序列域偏移明显等问题，构建了一套基于深度学习的无监督膝关节MRI分割与三维定量分析管线，
  并将其应用于前交叉韧带（ACL）撕裂与多发韧带膝关节损伤（MLKI）的影像定量比较中。
  \par 首先，在特征提取与跨域迁移方面，本文提出了一种基于 ResNet34 骨干、并融合实例归一化（IN）与空洞空间金字塔池化（ASPP）的改进型DANN网络。
  通过引入基于损失比率的动态冻结训练策略，尝试缓解对抗训练中的梯度失衡问题，并在目标域标签不参与训练的前提下，在临床FSE序列上获得了可用的骨与软骨自动分割结果（股骨 Dice达 $0.95$，软骨 Dice 达 $0.77$）。
  \par 其次，针对临床常规扫描层厚较大导致的各向异性断层问题，本文探索性地利用基于互信息的刚性配准结合深度图像先验（DIP）算法，
  评估了矢状、冠状、轴状多平面影像向 $0.5 mm$ 各向同性候选体积融合的可行性。
  该部分尚处于方法探索阶段，主要提供视觉连续性和自一致性重投影思路上的候选结果，尚不能作为经过定量验证的高分辨率重建模块。
  \par 最后，依托开发的欧氏距离三维厚度映射、kd-tree胫股关节影像近接触面积计算算法和BML候选高信号区域自动检测、粗分割与体积提取流程，本文对质量控制后纳入的 \num{579} 例患者进行了形态学提取与临床统计。
  结果显示，MLKI组相较于孤立ACL组具有更小的胫股关节影像近接触面积和更大的BML候选体积。
  这些发现提示两类损伤在关节空间关系与骨髓水肿范围上存在定量差异，但其与高能创伤“能量耗散”假说之间的因果关系仍需进一步研究验证。

\section{研究局限性}
  尽管本研究在无监督算法与临床应用结合上取得了一定进展，但仍存在部分局限性需要正视：
  \par 验证集规模有限：
  由于全手动标注三维软骨较为耗时，用以评估DANN模型性能的本地测试集仅包含了 30 例患者，
  其中用于第四章定量评价的是质控合格且与矢状面模型输入一致的150张FSE切片。
  未来需进一步扩大验证队列以提升模型泛化能力的统计学证据等级。
  \par 对照与消融实验不足：
  本文第四章主要报告最终改进型DANN模型在目标域测试集上的单模型结果，尚未系统完成 Source-only U-Net、标准DANN、DANN+IN、DANN+ASPP 和 DANN+动态冻结策略之间的定量对照。
  因此，当前结果可作为最终模型可用性的初步证据，但不能单独证明各结构改动的独立贡献。
  \par 三维分割的计算瓶颈：
  在多向高分辨率融合后，由于统计形状模型（SSM）拟合与 3D U-Net 高频子体积采样的计算开销较大，
  最终管线未能完成完整的三维端到端分割流程部署，在局部切片边界处仍偶有孤立假阳性预测。
  \par DIP重建验证不足：
  第五章的DIP多视图融合主要为探索性工作，当前缺少高分辨率金标准、重投影误差统计、层间平滑性定量指标以及下游分割性能对照。
  因此，DIP结果仅应作为候选三维输入和后续研究方向，而不应被解释为已经系统验证的三维重建成果。
  \par BML精细分割验证不足：
  本研究中的BML自动检测与分割主要指脂肪抑制FSE影像中高信号候选区域的自动检测、粗分割和体积提取，尚未建立大规模人工BML逐体素标注金标准，因此未报告BML病灶精细分割的Dice、ASD等准确性指标。
  相关结果应作为探索性影像量化指标解释，未来需要结合专家复核、外部验证数据和随访结局，进一步验证其病理对应关系和临床预测价值。

\section{未来展望}
  未来的工作可以从以下几个维度深化：
  \par 算法维度的时空拓展：
  尝试引入带有时间注意力机制的循环神经网络（Recurrent Neural Network, RNN）或基于图神经网络（Graph Neural Network, GNN）的拓扑约束层，
  在经过进一步定量验证的DIP候选重建体积或其他三维融合结果上进行全局分割，减少二维切片的局部视差。
  \par 纵向随访与预后建模：
  将本研究提取的形态学与影像组学特征（如 BML 候选体积、软骨厚度分布图）作为基线输入变量，结合患者随访期（1 年、2 年等）的临床量表，
  构建基于机器学习的预测模型，以刻画创伤后骨关节炎（PTOA）的纵向退变轨迹，并为韧带重建手术的干预时机提供更具前瞻性的人工智能（Artificial Intelligence, AI）辅助决策支持。
  
